\addchap{Abstract}

This thesis investigates automated verification of textual user-interface
requirements against ordered visual UI flows. The proposed output contract
combines a four-way requirement verdict with explicit screenshot evidence,
claim-level decisions, uncertainty reasons, and optional region-level
localization. The approach separates visible support from hidden system
properties and preserves uncertainty when the available visual states do not
justify a reliable positive or negative decision.

The empirical study uses 258 reviewed verification items across 13
Mind2Web-derived web flows. A controlled experiment varies automatic claim
decomposition and screenshot selection independently, complemented by model
comparisons, repeated runs, a flow-cluster bootstrap, an order-unavailable
robustness condition, and qualitative error analysis. Complete-flow prompting
provides the strongest label baseline in the controlled matrix. Restricting
raw requirements to lexically selected screenshots omits decisive states and
reduces label quality, while the effects of decomposition depend on the
available evidence. Explicit evidence representations remain useful for
traceability because they expose retrieval, interpretation, and aggregation
failures separately.

The evaluation uses a reviewed 13-flow benchmark and concerns what recorded
screenshots visibly establish. It does not infer hidden backend behavior, and
its results characterize this evidence regime rather than all UI designs and
implementations.
